\documentclass[11pt,aspectratio=169]{beamer}

\usetheme{Madrid}
\usecolortheme{seahorse}

\usepackage[utf8]{inputenc}
\usepackage{amsmath,amssymb}
\usepackage{xcolor}
\usepackage{booktabs}
\usepackage{tikz}
\usetikzlibrary{arrows.meta,positioning,fit,backgrounds}

\setbeamertemplate{navigation symbols}{}
\setbeamertemplate{footline}[frame number]

\title{Beyond the Syllabus}
\subtitle{Is Quadratic Really Necessary? Fine-Grained Lower Bounds and the Fall of
Bellman--Ford's 70-Year Barrier\\
\small Week 9 Supplement -- TOAE209/TOAH209: Design and Analysis of Algorithms}
\author{Foreign Trade University}
\date{}

\begin{document}

\begin{frame}
  \titlepage
\end{frame}

\begin{frame}{Outline}
  \tableofcontents
\end{frame}

% =================================================================
\section{Motivation}
% =================================================================

\begin{frame}{Why look beyond the syllabus?}
  \begin{itemize}
    \item This week you proved two upper bounds: edit distance / LCS in
    $O(nm)$ time via a 2-D DP table, and shortest paths with negative edges in
    $O(mn)$ time via Bellman--Ford.
    \item Two questions a curious student should ask next:
    \begin{enumerate}
      \item Is $O(nm)$ edit distance actually the \textbf{best possible}, or is a
      truly-subquadratic algorithm just waiting to be discovered?
      \item Dijkstra beats brute-force shortest paths when weights are non-negative.
      Is there an analogous ``fast'' algorithm once negative weights are allowed, or is
      $O(mn)$ Bellman--Ford as good as it gets?
      \item[$\to$] The first question was answered (no) in 2015. The second was answered
      (yes!) only in 2022 -- and sharpened again in \textbf{2026}.
    \end{enumerate}
  \end{itemize}
  \begin{alertblock}{How this supplement is different from a survey}
    Every result below is a specific, citable paper (author, venue, year, and an
    arXiv/DOI identifier), verified before this slide was written -- not a vague
    ``researchers have shown.''
  \end{alertblock}
\end{frame}

% =================================================================
\section{Is Quadratic Edit Distance Actually Necessary?}
% =================================================================

\begin{frame}{Recap: the bound you proved this week}
  \begin{block}{This week's lecture}
    Edit distance / sequence alignment between strings of length $n$ and $m$ is
    computable in $O(nm)$ time via a 2-D dynamic-programming table, using the
    ``what happens to the last character'' recurrence.
  \end{block}
  \begin{itemize}
    \item Every algorithms textbook presents $O(nm)$ as \emph{the} answer -- but is it
    optimal, or just the best anyone happened to find?
    \item Ruling out something faster needs a completely different kind of argument: a
    \textbf{conditional lower bound}, tying edit distance's difficulty to a
    widely-believed complexity-theoretic conjecture.
  \end{itemize}
\end{frame}

\begin{frame}{The conjecture: SETH}
  \begin{block}{Strong Exponential Time Hypothesis (SETH)}
    Boolean satisfiability on $n$ variables cannot be solved in
    $O(2^{(1-\varepsilon)n})$ time for \emph{any} fixed $\varepsilon > 0$ -- i.e.\ no
    algorithm beats exhaustive search by an exponentially-growing factor.
  \end{block}
  \begin{itemize}
    \item SETH is believed true (decades of SAT-solving research have found no
    counterexample) but is \emph{unproven} -- exactly like $P \ne NP$.
    \item A ``SETH-hardness'' result for problem $X$ says: \emph{if} $X$ has a
    truly-subquadratic algorithm, \emph{then} SETH is false. So a faster algorithm for
    $X$ would be a major breakthrough, not just an incremental improvement.
    \item The standard tool for proving this is a fine-grained reduction from
    \textbf{Orthogonal Vectors (OV)}: given two sets of $n$ boolean vectors in
    $\{0,1\}^{d}$ ($d = O(\log n)$), decide if some pair is orthogonal. OV is known to
    require $n^{2-o(1)}$ time under SETH.
  \end{itemize}
\end{frame}

\begin{frame}{The result: edit distance is SETH-hard}
  \begin{block}{Backurs \& Indyk, STOC 2015}
    ``Edit Distance Cannot Be Computed in Strongly Subquadratic Time (unless SETH is
    false).'' \emph{DOI:} \texttt{10.1145/2746539.2746612}.
  \end{block}
  \begin{itemize}
    \item They give a reduction turning an Orthogonal Vectors instance into two strings
    $X, Y$ of length $O(n)$ such that $\mathrm{edit\_distance}(X,Y)$ reveals whether an
    orthogonal pair exists.
    \item \textbf{Consequence:} an $O((nm)^{1-\varepsilon})$ edit-distance algorithm, for
    any $\varepsilon>0$, would give a truly-subquadratic OV algorithm, hence refute SETH.
    \item So the $O(nm)$ bound from this week's lecture is, modulo SETH, \textbf{exactly
    optimal} -- not a limitation of Kleinberg--Tardos's presentation, but a genuine wall.
  \end{itemize}
\end{frame}

\begin{frame}{What the lower bound does \emph{not} rule out}
  \begin{itemize}
    \item The Backurs--Indyk bound is about \textbf{exact}, \textbf{classical}
    algorithms only. Two escape hatches remain, and both are real, active research:
    \begin{enumerate}
      \item \textbf{Approximation.} Allowing a constant-factor approximation to the edit
      distance (rather than the exact value) sidesteps the reduction entirely.
      \item \textbf{Quantum computation.} A quantum algorithm might exploit structure the
      SETH reduction doesn't forbid.
    \end{enumerate}
    \item A 2018 result combines \emph{both} escape hatches at once.
  \end{itemize}
\end{frame}

\begin{frame}{The escape hatch: a quantum algorithm beats the wall}
  \begin{block}{Boroujeni, Ehsani, Ghodsi, HajiAghayi \& Seddighin, SODA 2018}
    ``Approximating Edit Distance in Truly Subquadratic Time: Quantum and MapReduce.''
    \emph{arXiv:} \texttt{1804.04178}.
  \end{block}
  \begin{itemize}
    \item They give a \textbf{quantum} algorithm running in $O(n^{1.858})$ time that
    approximates edit distance within a factor of $7$ -- genuinely subquadratic, and
    provably impossible \emph{classically} (exactly) under SETH.
    \item A second variant pushes the exponent down to $O(n^{1.781})$ for a larger
    approximation factor -- a direct time/accuracy trade-off.
    \item \textbf{Why this doesn't contradict Backurs--Indyk:} the SETH-hardness proof is
    for the \emph{exact} value; approximating within a constant factor is a
    genuinely easier problem, and quantum algorithms are a genuinely different
    computational model.
  \end{itemize}
\end{frame}

\begin{frame}{Section takeaway}
  \begin{enumerate}
    \item This week's $O(nm)$ DP for \emph{exact} edit distance is, modulo a
    widely-believed conjecture (SETH), the best any classical algorithm can ever do.
    \item That is a remarkably strong and clean statement -- most algorithms you meet in
    this course have \emph{no} matching lower bound at all.
    \item The wall is specifically about exact, classical computation: relax either
    constraint (allow approximation, allow quantum resources) and truly-subquadratic
    algorithms exist today.
  \end{enumerate}
\end{frame}

% =================================================================
\section{Breaking Bellman--Ford's 70-Year Barrier}
% =================================================================

\begin{frame}{Recap: the bound you proved this week}
  \begin{block}{This week's lecture}
    With negative edges allowed, Dijkstra's greedy strategy breaks (a settled vertex's
    distance can still decrease). Bellman--Ford instead relaxes every edge $n-1$ times,
    running in $O(mn)$ time, and can also detect negative cycles.
  \end{block}
  \begin{itemize}
    \item Bellman--Ford itself dates to the 1950s. For \emph{general} graphs with
    negative integer weights, no combinatorial algorithm improved meaningfully on this
    for decades -- unlike the non-negative case, where Dijkstra plus a good priority
    queue gives near-linear time immediately.
    \item Question: is $O(mn)$ actually necessary once negative weights are allowed, the
    way $O(nm)$ is (conditionally) necessary for edit distance above?
  \end{itemize}
\end{frame}

\begin{frame}{The 27-year plateau}
  \begin{block}{Goldberg, SIAM J.\ Computing, 1995 -- ``Scaling algorithms for the
  shortest paths problem''}
    A scaling technique (solve a coarser version of the weights, then refine) gives
    $O(m\sqrt{n}\,\log W)$ time, where $W$ bounds the largest negative weight magnitude.
  \end{block}
  \begin{itemize}
    \item This beats $O(mn)$ whenever $n$ is large relative to $\log W$, and it was the
    best \emph{combinatorial} bound for negative-weight SSSP for roughly 27 years.
    \item Near-linear time was known only for special cases (e.g.\ planar directed
    graphs) -- the general case seemed stuck.
    \item Contrast with this week's Dijkstra/Floyd--Warshall/DAG material: every one of
    those \emph{already} runs in close to $O(m+n)$ or $O(n^3)$ for all-pairs. Negative
    weights, for 70 years, cost you a factor of $\sqrt{n}$ or worse.
  \end{itemize}
\end{frame}

\begin{frame}{The 2022 breakthrough: near-linear time, at last}
  \begin{block}{Bernstein, Nanongkai \& Wulff-Nilsen, FOCS 2022 (pp.\ 600--611)}
    ``Negative-Weight Single-Source Shortest Paths in Near-linear Time.''
    \emph{arXiv:} \texttt{2203.03456}.
  \end{block}
  \begin{itemize}
    \item A \textbf{randomized} algorithm solving negative-weight SSSP in
    $\widetilde{O}(m\log W)$ time (near-linear, up to $\mathrm{polylog}$ factors) --
    matching, up to those factors, the non-negative-weight world Dijkstra already lives
    in.
    \item Core idea (very high level): combine the classical \emph{scaling} approach
    with \textbf{low-diameter graph decompositions} to repeatedly reduce the problem to
    a bounded-diameter instance that a cheap adaptation of Dijkstra can solve, without
    ever paying Bellman--Ford's full $n$ rounds of relaxation.
    \item This closed a gap that had stood, in some form, since Bellman--Ford itself.
  \end{itemize}
\end{frame}

\begin{frame}{The 2026 sequel: removing the randomness}
  \begin{block}{Haeupler, Jiang \& Saranurak, STOC 2026 -- ``Deterministic Negative-Weight
  Shortest Paths in Nearly Linear Time via Path Covers''}
    \emph{arXiv:} \texttt{2511.08551}.
  \end{block}
  \begin{itemize}
    \item Every near-linear-time algorithm since 2022 relied on \textbf{randomized}
    low-diameter decompositions. This 2026 paper gives the first fully
    \textbf{deterministic} nearly-linear algorithm, via a new structural primitive for
    directed graphs called a \textbf{path cover}.
    \item Running time: $\widetilde{O}(m)\cdot\log(nW)$ -- deterministic, no
    randomization anywhere, and it also detects negative cycles.
    \item \textbf{Small-world note:} Bernhard Haeupler is also the author (with
    Hershkowitz \& Li) behind the ``is Dijkstra universally optimal?'' result cited in
    Week 2's supplement -- the same research group is now rewriting the negative-weight
    side of shortest paths too.
  \end{itemize}
\end{frame}

\begin{frame}{The bound, over 70 years}
  \small
  \begin{center}
  \begin{tabular}{@{}lll@{}}
    \toprule
    Year & Bound (negative integer weights) & Source \\
    \midrule
    1950s & $O(mn)$ & Bellman--Ford \\
    1995 & $O(m\sqrt{n}\,\log W)$ & Goldberg (scaling) \\
    2022 & $\widetilde{O}(m\log W)$, \textbf{randomized} & Bernstein--Nanongkai--Wulff-Nilsen \\
    2026 & $\widetilde{O}(m)\log(nW)$, \textbf{deterministic} & Haeupler--Jiang--Saranurak \\
    \bottomrule
  \end{tabular}
  \end{center}
  \begin{exampleblock}{Perspective}
    Going from $O(mn)$ to near-linear is the same \emph{qualitative} leap as going from
    ``try every permutation'' to Karatsuba in Week 1 -- except here it took the research
    community 70 years, not one clever recursion.
  \end{exampleblock}
\end{frame}

\begin{frame}{Section takeaway}
  \begin{enumerate}
    \item $O(mn)$ Bellman--Ford is \textbf{not} optimal -- unlike edit distance, there
    was no lower-bound wall here, just an unsolved algorithmic question, and it fell.
    \item The fix came from importing a tool (low-diameter decompositions) from a
    different corner of algorithms research, not from a cleverer relaxation order.
    \item As of this supplement (Aug 2026), the deterministic version is only months
    old -- this is genuinely current research, not textbook history.
  \end{enumerate}
\end{frame}

% =================================================================
\section{Summary}
% =================================================================

\begin{frame}{Summary}
  \small
  \begin{enumerate}
    \item \textbf{Edit distance is provably hard:} Backurs \& Indyk (2015) showed
    $O(nm)$ is optimal for exact computation, modulo SETH -- one of the cleanest
    fine-grained lower bounds in the field.
    \item \textbf{But only for exact, classical computation:} Boroujeni et al.\ (2018)
    show a quantum algorithm can approximate edit distance in truly subquadratic time --
    the wall is real, but narrower than it first looks.
    \item \textbf{Bellman--Ford's $O(mn)$ bound, in contrast, was never a wall:}
    Bernstein--Nanongkai--Wulff-Nilsen (2022) reached near-linear time, and
    Haeupler--Jiang--Saranurak (2026) just removed the last randomization.
    \item \textbf{The general lesson:} ``textbook optimal'' always means
    \emph{optimal among known techniques at the time the textbook was written} --
    sometimes that turns out to be a provable wall, and sometimes it is simply the
    frontier of what anyone had tried yet.
  \end{enumerate}
  \begin{alertblock}{Looking ahead}
    Next week (Kleinberg \& Tardos, Ch.\ 7) moves to \textbf{network flow} -- where the
    story repeats once more: the textbook $O(VE^2)$ Edmonds--Karp bound is nowhere near
    the current frontier either.
  \end{alertblock}
\end{frame}

\end{document}
